\section{sprite/exec.hpp File Reference}
\label{exec_8hpp}\index{sprite/exec.hpp@{sprite/exec.hpp}}
Interface to the execution components. 

{\tt \#include \char`\"{}sprite/builtins.hpp\char`\"{}}\par
{\tt \#include \char`\"{}sprite/common.hpp\char`\"{}}\par
{\tt \#include \char`\"{}sprite/fingerprint.hpp\char`\"{}}\par
{\tt \#include \char`\"{}sprite/node.hpp\char`\"{}}\par
{\tt \#include \char`\"{}sprite/program.hpp\char`\"{}}\par
{\tt \#include \char`\"{}sprite/rewrite.hpp\char`\"{}}\par


Include dependency graph for exec.hpp:\begin{figure}[H]
\begin{center}
\leavevmode
\includegraphics[width=243pt]{exec_8hpp__incl}
\end{center}
\end{figure}


This graph shows which files directly or indirectly include this file:\begin{figure}[H]
\begin{center}
\leavevmode
\includegraphics[width=179pt]{exec_8hpp__dep__incl}
\end{center}
\end{figure}
\subsection*{Namespaces}
\begin{CompactItemize}
\item 
namespace {\bf sprite}
\end{CompactItemize}
\subsection*{Classes}
\begin{CompactItemize}
\item 
struct {\bf sprite::Yield\-Handler}
\begin{CompactList}\small\item\em Base of a class that handles computation results as they are generated. \item\end{CompactList}\item 
struct {\bf sprite::Yield\-Handler\_\-$<$ Output\-Iterator $>$}
\begin{CompactList}\small\item\em A concrete {\bf Yield\-Handler}{\rm (p.\,\pageref{structsprite_1_1YieldHandler})} for a particular output iterator type. \item\end{CompactList}\item 
struct {\bf sprite::Yield\-Handler\_\-$<$ void $>$}
\begin{CompactList}\small\item\em A yield handler that simply discards the output. \item\end{CompactList}\end{CompactItemize}
\subsection*{Enumerations}
\begin{CompactItemize}
\item 
enum {\bf Trace\-Option} \{ {\bf NO\_\-TRACE} = false, 
{\bf TRACE} = true
 \}
\end{CompactItemize}
\subsection*{Functions}
\begin{CompactItemize}
\item 
void {\bf execute} (Program const \&program, Node \&goal, Trace\-Option trace)\label{namespacesprite_a52}

\item 
template$<$typename Output\-Iterator$>$ void {\bf execute} (Program const \&pgm, Node \&goal, Output\-Iterator out, Trace\-Option trace=NO\_\-TRACE)\label{namespacesprite_a53}

\begin{CompactList}\small\item\em Alternate form of execute; results are written to the given iterator. \item\end{CompactList}\item 
bool {\bf is\_\-norm} (Node const \&node)\label{namespacesprite_a54}

\begin{CompactList}\small\item\em Traverses an expression and returns true if it is found to be in cnf. \item\end{CompactList}\item 
void {\bf pull\_\-tab} (Node \&g, Node\-Ptr const \&p)
\begin{CompactList}\small\item\em Performs the pull-tab transformation. \item\end{CompactList}\item 
void {\bf head\_\-normalize} (Node \&node)
\begin{CompactList}\small\item\em The head-normalizing (H) function from the fair scheme. \item\end{CompactList}\item 
void {\bf fair\_\-normalize} (Fingerprint const \&fp, Node \&node)\label{namespacesprite_a57}

\begin{CompactList}\small\item\em The normalizing (N) function from the fair scheme. \item\end{CompactList}\end{CompactItemize}


\subsection{Detailed Description}
Interface to the execution components. 



Definition in file {\bf exec.hpp}.

\subsection{Function Documentation}
\index{exec.hpp@{exec.hpp}!head_normalize@{head\_\-normalize}}
\index{head_normalize@{head\_\-normalize}!exec.hpp@{exec.hpp}}
\subsubsection{\setlength{\rightskip}{0pt plus 5cm}void head\_\-normalize (Node \& {\em node})\hspace{0.3cm}{\tt  [inline]}}\label{namespacesprite_file_a56}


The head-normalizing (H) function from the fair scheme. 

This is a generic H function that dispatches to the correct implementation for the given operation. 

Definition at line 135 of file exec.hpp.

References sprite::g\_\-program, sprite::Node::id(), sprite::Program::oper, and sprite::Node::tag().

Referenced by sprite::fair\_\-normalize(), and sprite::h\_\-func().

Here is the call graph for this function:\begin{figure}[H]
\begin{center}
\leavevmode
\includegraphics[width=140pt]{namespacesprite_a56_cgraph}
\end{center}
\end{figure}
\index{exec.hpp@{exec.hpp}!pull_tab@{pull\_\-tab}}
\index{pull_tab@{pull\_\-tab}!exec.hpp@{exec.hpp}}
\subsubsection{\setlength{\rightskip}{0pt plus 5cm}void pull\_\-tab (Node \& {\em g}, Node\-Ptr const \& {\em p})\hspace{0.3cm}{\tt  [inline]}}\label{namespacesprite_file_a55}


Performs the pull-tab transformation. 

\begin{Desc}
\item[Parameters:]
\begin{description}
\item[{\em g}]The source of the transformation. \item[{\em p}]The target of the transformation. Must be a direct descendant of g. \end{description}
\end{Desc}


Definition at line 118 of file exec.hpp.

References sprite::Node::clone(), sprite::Node::position(), and sprite::rewrite\_\-choice.

Referenced by sprite::fair\_\-normalize(), and sprite::h\_\-func().

Here is the call graph for this function:\begin{figure}[H]
\begin{center}
\leavevmode
\includegraphics[width=248pt]{namespacesprite_a55_cgraph}
\end{center}
\end{figure}
